\section{三価リンクを持つ matroidal 曲面の分類}
\label{trop:classification-section}

底空間の core を smooth tropical surface の点とみなせるかどうかには，
integral affine 構造を考える前から位相的な制約がある．
本節では二次元 matroidal fan の三価リンクを分類し，この制約を次数別に決定する．
特に，Petersen グラフが現れることの意味は，六次に特有の局所的な可能性として説明できる．

loopless matroid $M$ の Bergman fan を，対角方向を除いた convention
\[
 \mathcal B(M)\subset N_{M,\RR},\qquad
 N_M=\ZZ^{E(M)}/\ZZ(1,\ldots,1)
\]
で用いる．その次元は $\operatorname{rk}M-1$ である．
matroidal fan とその積を局所モデルにする意味での smooth tropical charts は
\cite[\S2.4]{Shaw} を参照されたい．
ここでいうリンクの同型型は，二価頂点の挿入・除去を許すグラフの位相型である．
従って，異なる fan subdivisions によって答えは変わらない．

\begin{theorem}[二次元 matroidal fan の三価リンク]
\label{trop:cubic-classification}
$M$ を loopless rank-three matroid とし，$\mathcal B(M)$ の原点のリンクに
一価または四価以上の点がないとする．parallel elements を簡約した後に $M$ が連結なら，
そのリンクと matroid は次の四つに限る．
\[
\begin{array}{c|c|c|c|c}
 |E(M)|&M&\text{三点 rank-two flats の数}&\text{リンク}&b_1\\\hline
 4&U_{3,4}&0&K_4&3\\
 5&M(K_4-e)&2&K_{3,3}&4\\
 6&M(K_4)&4&\text{Petersen}&6\\
 7&F_7&7&\text{Heawood}&8
\end{array}
\]
ここで $F_7$ は Fano matroid である．簡約後の ground set にラベルを付けたときの
各行の matroid の個数は，順に $1,15,30,30$ である．

非連結 matroids と Euclidean factors も許した任意の二次元 matroidal local model では，
上の四つに加えて円周と theta graph が現れる．特に，三価の分岐しか持たないリンクでは
\begin{equation}\label{trop:genus-bound}
 b_1\leq8
\end{equation}
である．この結論に matroid の体上での realizability は必要ない．
\end{theorem}

\begin{proof}
parallel elements $a,b$ は二点 circuit をなすので，Bergman 条件は $x_a=x_b$ を強制する．
各 parallel class から一元を選ぶ簡約は，支持集合の線形同相を誘導する．
従ってまず simple connected rank-three matroid に帰着する．

この場合，minimal nested set complex の頂点は，個々の元 $i$ と，
$|F|\geq3$ の rank-two flats $F$ である．辺は
\[
 i\text{---}F\quad(i\in F),\qquad
 i\text{---}j\quad\bigl(\operatorname{cl}\{i,j\}=\{i,j\}\bigr)
\]
である．これは，rank-two restriction が連結であるのが $|F|\geq3$ の場合であり，
二つの異なる rank-two flats の join が連結な全体 $E(M)$ になることから従う．
Feichtner--Sturmfels の triangulation theorem により，このグラフは Bergman link と
同じ位相型を持つ\cite[Theorem 4.1]{FeichtnerSturmfels}．

$F$ に対応する頂点の次数は $|F|$ である．従って，全ての rank-two flats は
二点または三点からなる．$n=|E(M)|$ とし，$i$ を通る三点 flats の数を $t_i$，
二点 flats の数を $s_i$ とすると
\begin{equation}\label{trop:point-count}
 2t_i+s_i=n-1,\qquad \deg(i)=t_i+s_i.
\end{equation}
異なる flats は $i$ 以外に共通点を持たないため，最初の式は各 $j\ne i$ を一回ずつ数える．
次数二は単なる細分点であり，それ以外の点は三価でなければならない．従って
\[
 n-1=2t_i+s_i\leq2(t_i+s_i)\leq6,
\]
すなわち $n\leq7$ である．また連結 rank-three matroid では $n\geq4$ である．

$n=4$ の場合，三点 flat があれば残る一点は coloop となるので，連結性から
$M=U_{3,4}$ に限る．その nested graph は $K_4$ である．

$n=5$ の場合，\eqref{trop:point-count} から各点が少なくとも一つの三点 flat に属する．
三点 flats は相異なるもの同士で高々一点を共有する．五点上では，二つの三点 flats は
ちょうど一点で交わり，三つ目は置けない．実際，最初の二つを $123,145$ とすると，
未使用の二点対は $24,25,34,35$ だけで，そこに三角形はない．
従って，この二つが全てである．共通点 $1$ の nested vertex は二価であり，
それを除くと $K_{3,3}$ になる．これは $M(K_4-e)$ である．

$n=6$ では $t_i=2,s_i=1$ が全ての $i$ について成り立つ．三点 flats は四つある．
これら四つを頂点とし，各元 $i$ が属する二つの flats を一本の辺で結ぶと，
四頂点の単純三価グラフを得る．単純性は二つの flats が二点を共有しないことによる．
従って，このグラフは $K_4$ であり，元の三点系はその四つの vertex stars，
すなわち Pasch configuration である．$K_4$ の各辺を反対辺に送る置換により，
この四つの stars は四つの三角形に移る．従って $M\cong M(K_4)$ である．
その nested graph は Petersen グラフである．
$K_{3,3}$ と Petersen の同定は，それぞれ
\cite[Examples 1.2 and 1.1]{FeichtnerSturmfels} にも明示されている．

$n=7$ では $t_i=3,s_i=0$ となり，全ての二点対がちょうど一つの三点 flat に含まれる．
これは七点の Steiner triple system である．点 $1$ を通る三つを
$123,145,167$ とし，残るラベルの交換で $246$ を選ぶと，
残りは $257,347,356$ に強制される．従って Fano plane が唯一であり，
その point--line incidence graph は Heawood graph である．
各グラフの $b_1$ は $|E|-|V|+1$ から表の通りになる．
ラベル付きの個数は，五点の場合 $5\cdot3=15$，六点の場合 $6!/24=30$，
七点の場合 $7!/168=30$ である．後二式の分母は，それぞれ $K_4$ の頂点置換群と
$\operatorname{GL}(3,\mathbb F_2)$ の位数である．

非連結 simple rank-three matroid は，連結 rank-two matroid と一つの rank-one component の
直和，または三つの rank-one components の直和である．前者は
\[
 \mathcal B(U_{2,k})\times\RR
\]
に同型で，リンクは $k$ 点集合の suspension，すなわち $K_{2,k}$ である．
三価性を満たす非平面の場合は $k=3$ だけで，二価頂点を除くと theta graph となる．
$k>3$ では suspension の両端が $k$ 価である．三つの rank-one components の場合は
$\RR^2$ であり，リンクは円周である．
明示的な $\RR$ factors を加えた二次元 charts も，coloops を加えた matroid として
同じ場合に含まれる．これで全ての二次元 matroidal local models を尽くす．
\end{proof}

\begin{remark}\label{trop:verification}
この分類は，有限次数の計算から一般の上界を推測したものではない．
\eqref{trop:point-count} が任意の ground set に対して $n\leq7$ を証明し，
残る四つの場合を上の議論が尽くしている．付属の整数検算は，$n=4,5,6,7$ の
全ての linear triple systems を backtracking で列挙し，basis exchange，連結性，
nested graph の同型，ラベル付き個数 $1,15,30,30$ を独立に検査する．
\end{remark}

\subsection{超曲面退化の次数別の障害}

\begin{theorem}[平面退化の core における局所位相分類]
\label{trop:degree-classification}
各 $T_{ij}$ を任意の marked pants tree とした底空間 $B(T)$ について，
各 core の近傍が二次元 matroidal fan の原点近傍に局所同相となる可能性は次で尽くされる．
\begin{enumerate}
\item $d=5$ では，全ての core link は $K_4$ であり，$\mathcal B(U_{3,4})$ によって
  局所位相型が実現される．
\item $d=6$ では，一つの core の局所位相型が matroidal であるための必要十分条件は，
  その link が Petersen graph であることである．全ての core についてこの条件を満たす
  ラベル付き分割系は正確に六個あり，$K_6$ の one-factorizations と対応する．
\item $d\geq7$ では，どの core の近傍も二次元 matroidal fan の原点近傍に局所同相でない．
  従って，core を保ったまま $B(T)$ を smooth tropical surface にすることはできない．
\end{enumerate}
第1，2項は局所位相型の実現を述べる．実際の fibration の action lattice と整合する
integral affine atlas の存在を含めない．
\end{theorem}

\begin{proof}
命題~\ref{topo:counts} より，全ての core link は三価で，
\[
 b_1(G_i)=\binom{d-2}{2}.
\]
fan の原点以外での三価性は，点の局所ホモロジーの階数からも判定できるため，
単なる局所同相を仮定しても定理~\ref{trop:cubic-classification} を適用できる．
より具体的には，punctured core germ の三価 locus の連結成分がリンクの頂点を，
その補集合の face germs と境界への incidence がリンクの辺を復元する．
従ってこの状況では，cone germs の局所同相はリンクのグラフ同相を誘導する．
$d=5$ のリンクは $K_4$，$d=6$ の第一 Betti 数は六であり，分類中では Petersen に限る．
その全 core での同時実現は定理~\ref{topo:petersen} によりちょうど六個である．
$d\geq7$ なら $\binom{d-2}{2}\geq10>8$ なので不可能である．
グラフの同相写像を cone に延長すれば，最初の二項の局所的十分性が従う．
\end{proof}

この障害は，退化成分を平面から高い次数の超曲面に替えても残る．
次では，異なる成分は transverse に交わり，三重点は transverse，四重点はないとする．
交点は個々に標識し，正の genus の曲線上でも，境界に extrema を持つ
simple-saddle Morse foliation を用いる．

\begin{proposition}[任意の degree partition に対する core link]
\label{trop:partition-obstruction}
$d=e_1+\cdots+e_k$，$e_i\geq1$，$k\geq3$ とし，次数 $e_i$ の滑らかな
超曲面 $S_i\subset\PP^3$ の union への退化を考える．
double curve $C_{ij}=S_i\cap S_j$ の genus と三重点数は
\begin{equation}\label{trop:double-curve-data}
 g_{ij}=1+\frac{e_i e_j(e_i+e_j-4)}2,
 \qquad n_{ij}=e_i e_j(d-e_i-e_j).
\end{equation}
その punctured curve の Reeb graph $T_{ij}$ は $b_1(T_{ij})=g_{ij}$ を満たし，
内部頂点数と内部辺数は
\[
 v_{ij}=2g_{ij}+n_{ij}-2=e_i e_j(d-4),\qquad
 a_{ij}=3g_{ij}+n_{ij}-3
\]
である．$S_i$ に対応する core link $\Gamma_i$ は連結三価グラフで，
\begin{equation}\label{trop:partition-core-genus}
 b_1(\Gamma_i)
 =\sum_{j\ne i}g_{ij}
  +\sum_{\substack{j<\ell\\j,\ell\ne i}}e_i e_j e_\ell-(k-1)+1
 =1+\frac{e_i(d-e_i)(d-4)}2.
\end{equation}
従って，$d\geq7$ では任意の degree partition と任意の marked Reeb graphs に対して
全ての core が matroidal local topology を持ち得ない．
$d=6$ でも，非線形成分が一つあれば，その core は matroidal でない．
従って六次で全 core が局所位相的に matroidal となる候補は全平面退化に限られ，
その分割系は前定理の六個で尽くされる．
\end{proposition}

\begin{proof}
adjunction と B\'ezout から \eqref{trop:double-curve-data} を得る．
$a,b,c\geq1$，$a+b+c=d$ に対して $abc\geq d-2$ なので
$n_{ij}\geq d-2\geq3$ である．曲線の pants decomposition は
$2g+n-2$ 個の pairs of pants を持ち，Reeb graph の内部辺数は
$3g+n-3$ となる．

$S_i$ の境界 divisor の dual multigraph は $k-1$ 頂点を持ち，
頂点 $j,\ell$ の間に $e_i e_j e_\ell$ 本の辺がある．
各頂点を $T_{ij}$ の内部部分に替えると，第一 Betti 数は
この dual graph の第一 Betti 数に $\sum_{j\ne i}g_{ij}$ を加えたものとなる．
これが \eqref{trop:partition-core-genus} の第一の式である．
$s=d-e_i=\sum_{j\ne i}e_j$ と書けば，その右辺は
\[
 1+\frac{e_i}{2}
 \left((e_i-4)s+\sum_{j\ne i}e_j^2+
                   2\sum_{\substack{j<\ell\\j,\ell\ne i}}e_j e_\ell\right)
 =1+\frac{e_i s(d-4)}2
\]
と整理できる．$1\leq e_i\leq d-2$ だから
\[
 e_i(d-e_i)\geq d-1,
 \qquad b_1(\Gamma_i)\geq1+\frac{(d-1)(d-4)}2=\binom{d-2}{2}.
\]
$d\geq7$ の主張は定理~\ref{trop:cubic-classification} に従う．
$d=6$ で $e_i=2,3,4$ の場合は $b_1=9,10,9$ であり，やはり上界八を超える．
\end{proof}

\section{五次の混合退化と matroidal core link の同時実現}
\label{q5:section}

次数だけから得られる第一 Betti 数は，異なる二重曲線が同じ core に入る方法を
記録しない．五次ではこの違いが本質的である．平面の core link の候補は $K_4$，
二次または三次成分の候補は $K_{3,3}$ であるが，それぞれを単独で実現できることから
全成分での同時実現は従わない．本節では，個々の三重点を標識した Reeb graphs を用いて，
五次の全 degree partitions に対するこの問題を決定する．

ここでいう \emph{matroidal} は，core の cone germ が二次元 Bergman fan の
原点近傍と局所同相であることを意味する．$K_4$ と $K_{3,3}$ を持つ fan は，
それぞれ $U_{3,4}$ と $M(K_4-e)$ の Bergman fan である
\cite[Examples 1.1--1.2 and Theorem 4.1]{FeichtnerSturmfels}．
以下は実際の Lagrangian fibration の action lattice による判定とは別の，
標識付きリンクの位相的な分類である．

\begin{theorem}[五次における全 core の同時分類]
\label{q5:classification}
$5=e_1+\cdots+e_k$，$e_i\geq1$，$k\geq3$ とし，
定理~\ref{part:global} の一般の degree-partition degeneration を考える．
各 double curve では，三重点を個々に標識し，境界に沿って円周葉を持ち，
内部には simple saddles のみを持つ Morse foliation を選ぶ．
全ての core links を同時に matroidal とする選択の存在は，成分の置換を除いて
次の表で尽くされる．
\[
\begin{array}{c|c|c|c}
 (e_1,\ldots,e_k)&\text{候補となる core links}
 &\text{全 core の同時実現}&\text{同時に matroidal となる最大個数}\\\hline
 (1,1,3)&K_4,K_4,K_{3,3}&\text{不可}&2\\
 (1,2,2)&K_4,K_{3,3},K_{3,3}&\text{不可}&2\\
 (1,1,1,2)&K_4,K_4,K_4,K_{3,3}&\text{可}&4\\
 (1,1,1,1,1)&K_4,K_4,K_4,K_4,K_4&\text{可}&5
\end{array}
\]
$(1,1,1,2)$ では，成分と個々の三重点のラベルを固定したとき，
全 core を matroidal とする葉ラベル付き抽象 Reeb graph の系はちょうど四個である．
これらは後述の parity 条件 \eqref{q5:odd-parity} により特徴付けられ，
各二点集合内の三重点ラベルの交換を許すと一つの軌道をなす．
全平面の場合の抽象 graph system は一つである．

ここで数えているものは，葉ラベル付き有限グラフの系である．
曲線への埋込みの marked isotopy class，葉に横断的な測度，
symplectic conjugacy class を数えてはいない．
各々の存在する graph system は，実際の標識付き double curves 上の
Morse foliations により実現でき，同じ曲線の両側で同一の foliation を使用できる．
\end{theorem}

\subsection{Divisor 成分を保つ graph expansion}

まず，曲線上の Morse data が core link の中にどのように残るかを記述する．
$T_a$ を，genus $g_a$，$n_a$ 個の境界円周を持つ曲面上の
simple-saddle Morse foliation の Reeb graph とし，$H_a$ を葉とそれに接する辺を
除いた内部部分とする．境界の各円周を一つの葉として扱う．
内部頂点が空でない場合，$H_a$ は連結であり，
\begin{equation}\label{q5:internal-counts}
 |V(H_a)|=2g_a+n_a-2,\qquad
 |E(H_a)|=3g_a+n_a-3,\qquad b_1(H_a)=g_a.
\end{equation}
一つの内部頂点からそれ自身に戻る辺は存在しない．実際，Morse function が
Reeb graph に誘導する実数値関数は各開辺上で単調なので，そのような辺の
両端の値を同じ critical value にすることはできない．
異なる二頂点を結ぶ多重辺はあり得る．

\begin{proposition}[誘導部分グラフとしての divisor blocks]
\label{q5:induced-blocks}
三重点のない SNC divisor の crossing graph を，各成分の $T_a$ で展開して得る
core link を $\Gamma$ とする．各 $T_a$ は内部頂点を少なくとも一つ持ち，
$T_a$ における内部頂点の価数は三とする．
このとき，$V(\Gamma)$ は $V(H_a)$ により分割され，次を満たす．
\begin{enumerate}
\item 各 $H_a$ は $\Gamma$ の連結な誘導部分グラフである．
\item 異なる blocks $H_a,H_b$ を結ぶ辺は，$D_a\cap D_b$ の個々の交点と
一対一に対応する．その端点は対応する標識付き葉が接していた内部頂点である．
\item $H_a$ から出る half-edges の数は $n_a$ であり，
\begin{equation}\label{q5:port-count}
 n_a=3|V(H_a)|-2|E(H_a)|.
\end{equation}
特に，$\Gamma$ が単純で girth $\ell$ を持ち，$g_a>0$ なら
\begin{equation}\label{q5:girth-obstruction}
 2g_a+n_a-2\geq\ell.
\end{equation}
\end{enumerate}
\end{proposition}

\begin{proof}
graph expansion は，各 $T_a$ の標識付き葉を，対応する crossing の相手の葉に
貼り，生じた二価点を除く操作である．この操作は内部頂点を同一視しない．
同じ $H_a$ の二頂点の間に新しい辺を作ることもない．従って各 block は誘導部分グラフで，
異なる blocks の間の辺は crossing labels そのものである．
三価性から各内部頂点の incident half-edges を数えれば
\eqref{q5:port-count} を得る．
$g_a=b_1(H_a)>0$ なら $H_a$ に cycle がある．単純グラフ $\Gamma$ の cycle は
少なくとも $\ell$ 個の相異なる頂点を含むので，\eqref{q5:internal-counts} から
\eqref{q5:girth-obstruction} が従う．
\end{proof}

この命題は，単なる $b_1(\Gamma)$ の比較より強い．例えば，genus one で
三つの punctures を持つ divisor は三頂点の cyclic block を強制するため，
全体が $K_{3,3}$ となることを妨げる．また，多重交点を一つの crossing type に
まとめてしまうと，第2項の情報を失い，以下の同時実現条件を検出できない．

\begin{lemma}[小さい matroidal links の block criterion]
\label{q5:block-criterion}
目標グラフ $G$ を $K_4$ または $K_{3,3}$ とする．
与えられた divisor incidence data が $G$ への graph expansion を持つための
必要十分条件は，$V(G)$ を二個以上の空でない connected induced blocks $V_a$ に分割し，
各 block の第一 Betti 数を $g_a$，出る half-edges の数を $n_a$ とし，
異なる blocks 間の辺を与えられた個々の crossing labels と対応させられることである．
この条件を満たす各 block の Reeb graph は，実際の genus $g_a$ の標識付き曲面上で実現できる．

候補となる proper connected induced blocks は次で尽くされる．
\[
\begin{array}{c|c|c|c}
G&H_a&g_a&n_a\\\hline
K_4&\text{一点}&0&3\\
K_4&K_{1,1}&0&4\\
K_4&K_3&1&3\\\hline
K_{3,3}&\text{一点}&0&3\\
K_{3,3}&K_{1,1}&0&4\\
K_{3,3}&K_{1,2}&0&5\\
K_{3,3}&K_{1,3}&0&6\\
K_{3,3}&K_{2,2}&1&4\\
K_{3,3}&K_{2,3}&2&3
\end{array}
\]
$K_{x,y}$ と $K_{y,x}$ は同じ行とする．
\end{lemma}

\begin{proof}
必要性は命題~\ref{q5:induced-blocks} による．$K_4$ の induced blocks は完全グラフ，
$K_{3,3}$ の connected induced blocks は一点または complete bipartite graph である．
後者について
\[
 b_1(K_{x,y})=(x-1)(y-1),\qquad n=3(x+y)-2xy
\]
なので，proper blocks を列挙すると表を得る．

十分性で注意すべき点は，正しい Euler characteristic を持つ抽象 graph が
必ず実数値 Morse function の Reeb graph になるわけではないことである．
必要な条件は，辺上で厳密に単調で，内部頂点に extrema を持たない高さ関数，
すなわち good orientation の存在である．この条件と Reeb graph の実現問題の
一般的な関係は \cite[Definition 2.2 and Theorem 5.6]{MichalakReeb} を参照されたい．
ここでは表中の各 graph に対して境界付きの構成を直接与える．

$H_a$ の各頂点 $v$ に $3-\deg_{H_a}(v)$ 本の葉を付けた graph を $T_a$ とする．
$H_a$ が tree の場合は，一つの葉を root とし，全ての辺を root から外向きに向ける．
各内部頂点は一つの入辺と二つの出辺を持つ．
$H_a=K_3$ または $K_{2,2}$ の場合は cycle の頂点に相異なる高さを付ける．
cycle 上で local minimum となる頂点には入る葉を，local maximum となる頂点には
出る葉を付ける．それ以外の頂点の葉はどちら向きでもよい．
各葉の高さを隣接頂点の適切な側に取れば，全ての辺で厳密に単調になる．
$H_a=K_{2,3}$ の場合は bipartition を $\{u_1,u_2\}$ と
$\{v_1,v_2,v_3\}$ とし，
\[
 h(v_1)<h(u_1)<h(v_2)<h(u_2)<h(v_3)
\]
とする．三つの葉は $v_1,v_2,v_3$ に付き，$v_1$ の葉を入辺，
$v_3$ の葉を出辺とすればよい．$v_2$ の葉の向きは任意である．
これで表中の全ての $T_a$ に good orientation を与えた．

各三価頂点を一つの saddle を持つ oriented pair of pants に置き換え，
高さに従って円筒を貼る．入辺が一つの場合は一円周が二円周へ分かれる saddle，
入辺が二つの場合はその逆を使う．相異なる内部頂点の critical values は全て異なるように取る．
残った葉の端には radial collar を付ける．この構成の曲面は連結かつ orientable で，
境界数は $n_a$，Euler characteristic は $-|V(H_a)|=2-2g_a-n_a$ である．
従って genus は $g_a$ である．標識付き境界を保つ orientation-preserving
diffeomorphism により，与えられた実際の punctured curve へ移すことができる．
この diffeomorphism は境界 collar の指定された radial coordinates と整合するように選べる．
それぞれの crossing label で葉を貼れば，元の $G$ が復元される．
\end{proof}

この十分性は，表中の blocks について確認した good orientation を用いている．
任意の target graph について，genus と incidence の数だけを同じにすれば
Morse 実現が自動的に従う，という主張ではない．

\subsection{五次の数値と三重点ラベル}

以下では，成分を $A,B,C,Q$ などで表し，double curve $C_{ij}$ に対応する内部 graph を
$H_{ij}$ と書く．五次では \eqref{trop:double-curve-data} が
\begin{equation}\label{q5:curve-data}
 g_{ij}=1+\frac{e_i e_j(e_i+e_j-4)}2,
 \quad n_{ij}=e_i e_j(5-e_i-e_j),
 \quad |V(H_{ij})|=e_i e_j
\end{equation}
となる．また命題~\ref{trop:partition-obstruction} より
\begin{equation}\label{q5:core-genus}
 b_1(\Gamma_i)=1+\frac{e_i(5-e_i)}2.
\end{equation}
従って，$e_i=1$ の core は四頂点，$e_i=2,3$ の core は六頂点を持つ．
定理~\ref{trop:cubic-classification} によれば，それらの matroidal 候補は
それぞれ $K_4$ と $K_{3,3}$ だけである．

四つの葉を持つ genus-zero Reeb graph は，二つの三価頂点を一本の内部辺で結んだ
quartet である．四点集合 $P$ 上の葉ラベル付き quartet は，$P$ の unordered な
$2+2$ 分割と一対一に対応する．この三つの分割の集合を $\mathcal Q(P)$ と書く．
ここで $P$ の元は triple \emph{type} ではなく，実際の個々の三重点である．

\subsection{二平面と三次成分：\texorpdfstring{$(1,1,3)$}{(1,1,3)}}

$A,B$ を平面，$C$ を三次成分とする．共通の三重点集合 $P=A\cap B\cap C$ は三点である．
$C_{AB}$ は genus zero，三 punctures を持つので，$H_{AB}$ は一点である．
$C_{AC},C_{BC}$ はそれぞれ genus one，三 punctures を持ち，
$H_{AC},H_{BC}$ は三頂点三辺の connected graphs である．

\begin{proposition}\label{q5:113}
$(1,1,3)$ では三次成分の core は，どの marked Morse foliations を選んでも
matroidal にならない．両方の平面 core を $K_4$ にすることはできる．
その場合，三次成分の core link は triangular prism である．
\end{proposition}

\begin{proof}
三次成分の link には $H_{AC}$ と $H_{BC}$ が induced blocks として入る．
もし link が $K_{3,3}$ なら，genus one，三頂点の block を含むことになるが，
これは補題~\ref{q5:block-criterion} の表にない．直接には，三頂点三辺の単純 graph は
三角形であり，bipartite graph に入らない．$H_{AC}$ または $H_{BC}$ が非単純なら，
やはり単純な $K_{3,3}$ の部分グラフにならない．

一方，平面 $A$ の link が $K_4$ となるための必要十分条件は，
$H_{AC}$ が三角形であり，各頂点に $P$ の葉が一つずつ付くことである．
$H_{AB}$ の一点は三本の crossing edges によってその三角形の各頂点に結ばれる．
平面 $B$ についても同様である．三角形に三つの葉を付けた Reeb graph は
補題~\ref{q5:block-criterion} により actual elliptic curve 上で実現できる．
両方をこのように取ると，三次成分の link は二つの三角形を三本の対応する
$P$-labelled edges で結んだもの，すなわち triangular prism である．
\end{proof}

\subsection{一平面と二つの二次成分：\texorpdfstring{$(1,2,2)$}{(1,2,2)}}

$A$ を平面，$B,C$ を二次成分とし，四つの共通三重点の集合を $P$ とする．
$C_{AB}$ と $C_{AC}$ は四 punctures を持つ rational curves なので，
その quartets を
\[
 \sigma_{AB},\sigma_{AC}\in\mathcal Q(P)
\]
で表す．$C_{BC}$ は四 punctures を持つ genus-one curve であり，
$H_{BC}$ は四頂点四辺の connected graph である．

\begin{proposition}\label{q5:122}
$(1,2,2)$ では，全ての core を同時に matroidal にすることはできない．
ただし，三つの core のうち任意に指定した二つを同時に matroidal にできる．

より正確には，少なくとも一つの二次成分の core が $K_{3,3}$ なら，
$H_{BC}$ は四角形であり，その各頂点に一つの葉が付く．
四角形の bipartition が定める $P$ の $2+2$ 分割を $\tau$ とすると，
\begin{align}
 \Gamma_A\cong K_4&\quad\Longleftrightarrow\quad
      \sigma_{AB}\ne\sigma_{AC},\label{q5:122-plane}\\
 \Gamma_B\cong K_{3,3}&\quad\Longleftrightarrow\quad
      \sigma_{AB}=\tau,\label{q5:122-B}\\
 \Gamma_C\cong K_{3,3}&\quad\Longleftrightarrow\quad
      \sigma_{AC}=\tau.\label{q5:122-C}
\end{align}
\end{proposition}

\begin{proof}
$K_{3,3}$ の induced block のうち四頂点で genus one を持つものは $K_{2,2}$ に限る．
その各頂点の内部次数は二なので，残りの一つの half-edge が $P$ の葉となる．
従って，quadric core が matroidal なら上記の四角形型に限られる．

plane core では，$H_{AB}$ と $H_{AC}$ はそれぞれ一本の辺である．
二つの端点間を結ぶ crossing edges の本数は，対応する二つの split parts の
共通部分の大きさに等しい．二つの $2+2$ 分割が異なれば，全ての共通部分は一点である．
従って四本の crossing edges は $K_{2,2}$ をなし，二本の内部辺を加えると $K_4$ になる．
分割が同じなら二組の二重辺を生じる．これが \eqref{q5:122-plane} である．

quadric $B$ の link は，四角形 $H_{BC}$ に，辺 $H_{AB}$ の両端を加えた graph である．
$K_{3,3}$ となるためには，新しい二頂点がその bipartition の異なる側に属し，
それぞれ四角形の反対側の二頂点に接続しなければならない．
これは $\sigma_{AB}=\tau$ と同値で，そのとき得られる三価 bipartite graph は
実際に $K_{3,3}$ である．$C$ についても同じ議論を適用する．

両 quadric cores を $K_{3,3}$ にすると
$\sigma_{AB}=\sigma_{AC}=\tau$ が強制され，plane core の条件と矛盾する．
これで同時実現の不可能性が示された．両 quadrics だけを実現するには三つの分割を全て同じにし，
plane と $B$ を実現するには $\sigma_{AB}=\tau\ne\sigma_{AC}$ とする．
plane と $C$ の場合も同様である．四角形型と quartet 型は
補題~\ref{q5:block-criterion} により実際の marked curves 上で実現できる．
\end{proof}

\subsection{三平面と一つの二次成分：\texorpdfstring{$(1,1,1,2)$}{(1,1,1,2)}}

平面を $A,B,C$，二次成分を $Q$ とする．三平面の共通点を $q$ とし，
\[
 P_{AB}=A\cap B\cap Q,\quad
 P_{AC}=A\cap C\cap Q,\quad
 P_{BC}=B\cap C\cap Q
\]
と書く．これらは相異なる二点集合である．plane--plane curves 上の Reeb graphs は
全て三葉の $Y$-graph である．plane--quadric curves 上の quartets は
\[
 \sigma_A\in\mathcal Q(P_{AB}\sqcup P_{AC}),\quad
 \sigma_B\in\mathcal Q(P_{AB}\sqcup P_{BC}),\quad
 \sigma_C\in\mathcal Q(P_{AC}\sqcup P_{BC})
\]
で定まる．

各 quartet に一つずつ存在する，二つの二点集合をそのまま分離する split を
\emph{grouped} と呼ぶ．他の二つは，各側が両方の二点集合から一点ずつを含むので
\emph{crossed} と呼ぶ．

\begin{proposition}[三つの quartet の parity]
\label{q5:1112}
各 plane core が $K_4$ となるための必要十分条件は，その plane に対応する
quartet が crossed であることである．$Q$ の core が $K_{3,3}$ となるための
必要十分条件は，三つの quartets が全て crossed で，以下の parity が odd となることである．

$P_{XY}=\{p_{XY}^0,p_{XY}^1\}$ と各二点集合を列挙する．
crossed splits を，$s\in\mathbb F_2$ に対して
\begin{align*}
 \sigma_A &: \{p_{AB}^s,p_{AC}^{s+\epsilon_A}\},\\
 \sigma_B &: \{p_{AB}^s,p_{BC}^{s+\epsilon_B}\},\\
 \sigma_C &: \{p_{AC}^s,p_{BC}^{s+\epsilon_C}\}
\end{align*}
という二つの parts で表す．このとき必要十分条件は
\begin{equation}\label{q5:odd-parity}
 \epsilon_A+\epsilon_B+\epsilon_C=1\quad\text{in }\mathbb F_2
\end{equation}
である．全て crossed で parity が even の場合，$Q$ の link は triangular prism となる．
\end{proposition}

\begin{proof}
plane $A$ の link は，$H_{AB}$，$H_{AC}$ の二つの一点 blocks と，
quartet の内部辺 $H_{AQ}$ からなる．前二者は $q$ による一本の辺で結ばれる．
$P_{AB}$ と $P_{AC}$ は，それぞれ前二者を quartet の端点へ結ぶ二本の辺である．
split が crossed なら，これらは四つの可能な端点対を一回ずつ結ぶので $K_4$ を得る．
grouped なら二重辺を生じる．$B,C$ についても同様である．

まず quadric core が $K_{3,3}$ なら，三つの quartets は全て crossed でなければならない．
例えば $\sigma_A$ が grouped とする．$P_{AB}$ の二本の crossing edges は
$H_{AQ}$ の同じ端点から出る．$\sigma_B$ も grouped なら，その二辺は
$H_{BQ}$ の同じ端点に入り，二重辺となる．$\sigma_B$ が crossed なら，
$H_{BQ}$ の異なる端点に入り，その内部辺と合わせて三角形を作る．
いずれも単純 bipartite graph $K_{3,3}$ に反する．

そこで全て crossed とする．$H_{AQ},H_{BQ},H_{CQ}$ の端点を，
上の split parts に対応して $A_s,B_s,C_s$ と書く．
三本の内部辺は
\[
 A_0A_1,\qquad B_0B_1,\qquad C_0C_1
\]
であり，六本の crossing edges は
\begin{equation}\label{q5:voltage-edges}
 A_sB_s,\qquad A_sC_{s+\epsilon_A},\qquad
 B_sC_{s+\epsilon_B+\epsilon_C}
 \qquad(s\in\mathbb F_2)
\end{equation}
である．従って crossing subgraph は三角形の二重被覆をなし，一周したときの
sheet change は $\epsilon_A+\epsilon_B+\epsilon_C$ に等しい．

parity が even なら被覆は二つの三角形の disjoint union である．
各内部辺はその二つの対応する頂点を結ぶので，graph は triangular prism である．
parity が odd なら被覆は一本の六角形となる．その deck involution は
六角形の対頂点を交換し，三本の内部辺はその対頂点を結ぶ．
得られる graph は $K_{3,3}$ である．
\end{proof}

\begin{corollary}[四つの標識付き系と完全な quartet 数え上げ]
\label{q5:four-systems}
$(1,1,1,2)$ における $3^3=27$ 個の quartet systems のうち，
matroidal cores の個数が $4,3,2,1,0$ となるものは，それぞれ
\[
 4,\quad4,\quad12,\quad6,\quad1
\]
個である．全 core を実現する四つは三つの二点集合内のラベル交換に関して一つの軌道をなす．
\end{corollary}

\begin{proof}
全て crossed の八個のうち，odd parity の四個が全 core を実現し，
even parity の四個は三つの plane cores だけを実現する．
少なくとも一つが grouped なら quadric core は matroidal でないので，
matroidal cores の個数は crossed quartets の個数に一致する．
grouped が一つ，二つ，三つの系はそれぞれ
$\binom31 2^2=12$，$\binom32 2=6$，$1$ 個である．

$P_{AB}$ 内のラベルを交換すると $\epsilon_A,\epsilon_B$ が同時に変わり，
$P_{AC}$ なら $\epsilon_A,\epsilon_C$，$P_{BC}$ なら
$\epsilon_B,\epsilon_C$ が同時に変わる．従って parity は列挙の選び方に依存しない．
これらの二 bit の反転は，odd parity を持つ四つの vectors に推移的に作用する．
\end{proof}

例えば $(\epsilon_A,\epsilon_B,\epsilon_C)=(1,0,0)$ とすれば，
crossing subgraph の六角形は
\[
 A_0-B_0-C_0-A_1-B_1-C_1-A_0
\]
であり，残る三辺が対頂点を結ぶ．この記述は，実際の七個の三重点
$q,P_{AB},P_{AC},P_{BC}$ のラベルを全て保持している．

\subsection{Actual marked foliations と主定理の結論}

以上の存在部分で使った内部 graphs は，一点，辺，三角形，四角形だけである．
補題~\ref{q5:block-criterion} の構成は，それぞれを指定された punctures を持つ
実際の rational または genus-one curve 上で実現する．
特に，genus-one graph を rational pants tree に置き換えたり，
同じ triple type の複数の点を一つにまとめたりしていない．

各 $C_{ij}$ 上でこのような foliation を一度だけ選び，両側の成分に同じものを指定する．
各 puncture の十分小さい固定円板では radial foliation とし，
crossing の product germ を共通に選ぶ．その後，定理~\ref{snc:main} の相対構成と
定理~\ref{part:global} の接着を適用する．
これらの構成は指定した曲線上の foliation を保持するので，
ここで決定した graph systems は実際の fibrations の core links として現れる．

なお，この実現は指定した抽象型を持つ marked foliation の\emph{存在}を与える．
同じ leaf-labelled Reeb graph を持つ任意の二つの foliations が
標識を固定して isotopic，あるいは symplectically conjugate であるとは主張しない．
後者には measured Reeb graph の情報も必要である
\cite{DMT,IKM}．また，Bergman fan の abstract cone と同相であることから，
action-compatible integral affine chart や，所定の symplectic filling との同一性は従わない．

命題~\ref{q5:113}，\ref{q5:122}，\ref{q5:1112} と
系~\ref{q5:four-systems} が三つの mixed partitions を尽くす．
残る全平面 partition では，どの $C_{ij}$ も三 punctures を持つ rational curve で，
その内部 graph は一点である．各 core では四つの一点 blocks が全て一回ずつ結ばれ，
$K_4$ を得る．これで定理~\ref{q5:classification} が従う．



\section{Petersen star の整数格子と torus cycles}
\label{trop:petersen-lattice-section}

\subsection{標準 fan の primitive rays と balancing}

$K_4$ の辺集合を
$E=\{12,13,14,23,24,34\}$ とし，
\[
 N=\ZZ^E/\ZZ\mathbf 1
\]
とする．この格子の rank は五であり，fan 自体の次元二とは異なる．
Petersen graph を，五点集合の二元部分集合を頂点とし，互いに素な部分集合を結ぶ
Kneser graph $KG(5,2)$ として表示する．その十本の rays を
\begin{equation}\label{trop:ray-formula}
 r_{ab}=[e_{ab}]\quad(1\leq a<b\leq4),\qquad
 r_{a5}=\left[\sum_{\substack{b<c\\b,c\in\{1,2,3,4\}\setminus\{a\}}}e_{bc}\right]
 \quad(1\leq a\leq4)
\end{equation}
で定める．二次元 cones は $S\cap T=\varnothing$ を満たす
$\RR_{\geq0}r_S+\RR_{\geq0}r_T$ の十五個である．
これは $\mathcal B(M(K_4))$ の minimal nested fan であり，
$M_{0,5}^{\mathrm{trop}}$ の標準 fan model と同型である
\cite[\S4]{ArdilaKlivans}\cite[arXiv版, Example 2.28]{Shaw}．

基底を $[e_{12}],[e_{13}],[e_{14}],[e_{23}],[e_{24}]$ とすると，座標は次の通りである．
\[
\begin{array}{c|rrrrr}
 S&r_S^{(1)}&r_S^{(2)}&r_S^{(3)}&r_S^{(4)}&r_S^{(5)}\\\hline
 12&1&0&0&0&0\\
 13&0&1&0&0&0\\
 14&0&0&1&0&0\\
 23&0&0&0&1&0\\
 24&0&0&0&0&1\\
 34&-1&-1&-1&-1&-1\\
 15&-1&-1&-1&0&0\\
 25&-1&0&0&-1&-1\\
 35&1&0&1&0&1\\
 45&1&1&0&1&0
\end{array}
\]

\begin{proposition}\label{trop:petersen-unimodular}
上の rays は primitive であり，各二次元 cone の二本の generators は，
その実線形 span と $N$ との交わりの格子基底をなす．
また，重みを一とした balancing identity は
\begin{equation}\label{trop:petersen-balancing}
 \sum_{T\cap S=\varnothing}r_T=r_S
\end{equation}
である．Petersen adjacency matrix を $A$ とすると，
\begin{equation}\label{trop:petersen-presentation}
 N\cong\operatorname{coker}(A-I:\ZZ^{10}\longrightarrow\ZZ^{10})\cong\ZZ^5,
\end{equation}
ここで十個の generators はそれぞれ $r_S$ に送られる．
\end{proposition}

\begin{proof}
非空な真部分集合の indicator vector の対角方向を除いた像は primitive である．
二本とも singleton-edge rays である場合，それらの座標と別の一座標を選ぶと，
二本の vectors と $\mathbf1$ の三次 minor に $\pm1$ がある．
一方が triangle-flat ray の場合，もう一方の辺はその三角形に含まれる．
その辺，同じ三角形の別の辺，三角形外の辺を選ぶと，やはり三次 minor が $\pm1$ になる．
これは $N$ における rank-two sublattice の saturation を示す．
二本とも triangle-flat rays である cone は存在しない．

\eqref{trop:petersen-balancing} は \eqref{trop:ray-formula} に代入すれば従う．
例えば $S=a5$ では隣接する三本の singleton-edge rays の和そのものであり，
$S=12$ では $r_{34}+r_{35}+r_{45}=r_{12}$ となる．
$a5$ に対応する四つの関係を使うと，$r_{a5}$ は残る六個の edge generators の和で消去できる．
残りの六つの関係はいずれも $\sum_{ab\in E}e_{ab}=0$ に帰着する．
従って，これは自由格子 $\ZZ^6/\ZZ\mathbf1$ の整数係数 presentation であり，torsion はない．
\end{proof}

標準的な edge length の整数単位との整合も確認できる．
Gathmann--Kerber--Markwig\cite[Construction 3.6, Theorem 3.7]{GKM}は，
長さ一の internal edge を持つ split tree の vectors $v_S$ の整数 span を
$M_{0,5}^{\mathrm{trop}}$ の lattice $\Lambda_5$ とする．
これらは \eqref{trop:petersen-balancing} と同じ関係を満たし，$\Lambda_5$ の rank は五である．
従って \eqref{trop:petersen-presentation} から，$r_S\mapsto v_S$ は格子同型となる．
以下ではこの unit-edge lattice の正規化を用いる．

\subsection{六つの stars への同時ラベル付け}

定理~\ref{topo:petersen} の one-factorization を
$\{\infty\}\cup\mathbb F_5$ 上に取り，
$M_c=\{\infty c,(c+1)(c-1),(c+2)(c-2)\}$ とする．
$i\in\mathbb F_5$ に対して
\[
 f_i(x)=\frac1{x-i},\qquad f_i(i)=\infty,\qquad f_i(\infty)=0,
 \qquad f_\infty=\operatorname{id}
\]
と置く．translation と inversion がこの factorization を保つため，
$f_i$ は $i$ の star を $\infty$ の star に写す．
後者の頂点を $A_c,B_c$ とした表示に対し
\begin{equation}\label{trop:kneser-labels}
 A_c\longmapsto\{2c,2c+1\},\qquad
 B_c\longmapsto\{2c+2,2c+4\}\quad\text{in }\mathbb F_5
\end{equation}
と置くと，隣接は二元部分集合の非交差性と一致する．
最後に $0,1,2,3,4$ を $1,2,3,4,5$ と読み替え，
\eqref{trop:ray-formula} を割り当てれば，全六 stars について primitive rays を明示できる．
他の五つの factorization もラベルの置換で得られる．

この構成は各 core の抽象的な fan lattice を指定する．
三角形・二角形の内部を通る lattice transition maps は，このラベル付けのデータには含まれない．
また，$N$ は fan の ambient tangent lattice であり，Lagrangian torus の
$H_1(\TT^2;\ZZ)$ が action cotangent lattice を与えることとは convention が異なる．
両者を同一視する場合には，face ごとの rank-two lattices と dualization を明示する必要がある．

\subsection{common-circle invariant と divisor incidence の復元}

SNC divisor に付随する contact-link fibration $S\to\Gamma$ を考える．
各三価頂点 $v$ の小さい閉近傍の逆像 $Q_v$ は，pair of pants $P$ 上の circle bundle であり，
$H^2(P;\ZZ)=0$ によって $S^1\times P$ と同相である．
三つの boundary tori の $H_1$ は $H_1(Q_v;\ZZ)\cong\ZZ^3$ に injective に入り，
その三つの像の共通部分は primitive rank-one subgroup である．その generator を
符号を除いて $\alpha_v$ と書く．

\begin{definition}\label{trop:delta-definition}
辺 $e=vw$ 上の regular torus の homology を辺に沿って同定し，
両端の common-circle generators を同じ $H_1(\TT^2;\ZZ)$ に移す．
\[
 \delta(e)=|\det(\alpha_v,\alpha_w)|
\]
を \emph{common-circle determinant} と呼ぶ．これは基底，向き，trivialization に依存しない．
\end{definition}

\begin{proposition}[cycle data からの divisor 復元]
\label{trop:divisor-recovery}
$D=\bigcup_aD_a$ を SNC divisor とし，各 punctured component が
$n_a\geq3$ 個の punctures を持つとする．genus $g_a$ の成分には
$2g_a+n_a-2$ 個の simple saddles を持つ marked pants model を選ぶ．
その link fibration について
\[
 \delta(e)=
 \begin{cases}
 0,&\text{同一の }D_a\text{ の Reeb graph 内の辺},\\
 1,&\text{SNC crossing に対応する辺}
 \end{cases}
\]
である．全ての頂点を含む $\delta=0$ 部分グラフの連結成分は，各 $D_a$ の内部 Reeb graph であり，
その第一 Betti 数は $g_a$ である．各連結成分を一点に縮約すると，元の dual multigraph を復元できる．
従って fiber-preserving な link-fibration germ は，dual incidence と各成分の genus を記憶する．
有理成分だけの場合，$\delta=0$ の連結成分は内部 pants trees である．
\end{proposition}

\begin{proof}
$Q_v\cong S^1\times P$ の三つの tori の homology images を
$\langle\alpha,\beta_1\rangle$，$\langle\alpha,\beta_2\rangle$，
$\langle\alpha,\beta_1+\beta_2\rangle$ と書けば，その共通部分は $\ZZ\alpha$ である．
この $\alpha$ は成分 $D_a$ の normal circle に一致する．
同じ成分内の annulus を通ると，両端の normal circles は同じ cycle なので determinant は零である．
crossing の product chart では，両成分の normal circles が torus の整数基底になるので determinant は一である．
内部 Reeb graph は連結であり，葉を除いても第一 Betti 数 $g_a$ は変わらない．
従って，零 determinant の連結成分を潰せば，残る一本一本の辺が正確に SNC crossing となる．
これらの構成は全て torus homology と vertex-block inclusions から定義されるため，
fiber-preserving 同値で保たれる．
\end{proof}

\begin{remark}
saddle を持たない二価の有理成分は，この復元定理の仮定に含めない．
そのような成分を許した一般的な boundary homology の完全列とは，適用範囲が異なる．
\end{remark}

$m$ 本の一般位置の直線の場合，各成分は $m-1$ 個の punctures を持つので
\[
 \#\{e:\delta(e)=0\}=m(m-4),\qquad
 \#\{e:\delta(e)=1\}=\binom m2.
\]
零辺の縮約は $K_m$ を与える．特に $m=5$ の Petersen model では，
五本の internal matching edges が $\delta=0$，残る十本が $\delta=1$ である．

一方，$\overline M_{0,5}$ は $\PP^2$ の一般位置の四点での blow-up であり，
その境界は十本の $(-1)$-curves からなる\cite[Theorem 4.3.3]{KapranovChow}．
四本の exceptional curves と六本の二点を結ぶ直線の strict transforms を用いれば，
その交点グラフは Petersen で，各成分は三点で他成分と交わることが直接分かる．
従って，この compactification の標準 pants link では十五辺の全てで $\delta=1$ となる．

\begin{corollary}\label{trop:not-standard-M05}
五直線補集合 $W_5$ に由来する Petersen link fibration は，
$\overline M_{0,5}\setminus M_{0,5}$ に由来する標準 SNC link fibration と
fiber-preserving 同値ではない．さらに，両者の boundary three-manifolds は同相でない．
\end{corollary}

\begin{proof}
最初の主張は $\delta=0$ の辺数の差から従う．第二の主張は整数ホモロジーで区別できる．
有理 SNC divisor の plumbing neighborhood $N_D$，境界 $S_D$，
intersection matrix $Q_D$，dual graph $G_D$ に対して
\begin{equation}\label{trop:boundary-homology}
 0\longrightarrow\operatorname{coker}(Q_D:\ZZ^r\to\ZZ^r)
 \longrightarrow H_1(S_D;\ZZ)
 \longrightarrow H_1(G_D;\ZZ)\longrightarrow0
\end{equation}
が成り立つ．実際，$N_D$ は $D$ に retract し，$H_2(N_D)=\ZZ^r$，
$H_1(N_D)=H_1(G_D)$ は自由群，$H_3(N_D)=0$ である．
Poincar\'e--Lefschetz duality と対の完全列により，左の写像は交点形式 $Q_D$ となる．

五直線の場合は $Q_D=J_5$，従って cokernel は $\ZZ^4$，$b_1(K_5)=6$ であり，
\[
 H_1(S_{W_5};\ZZ)\cong\ZZ^{10}.
\]
標準 $\overline M_{0,5}$ 境界では $Q_D=A-I$ であり，
\eqref{trop:petersen-presentation} と $b_1(\text{Petersen})=6$ により
\[
 H_1(S_{M_{0,5}};\ZZ)\cong\ZZ^{11}.
\]
\eqref{trop:boundary-homology} の右端が自由群なので，いずれも extension は分裂する．
\end{proof}

\begin{remark}\label{trop:standard-model-scope}
この結果が禁じるのは，標準 $M_{0,5}$ の SNC contact model と整合する局所同定である．
抽象的な Bergman fan atlas だけからその plumbing filling が強制されるとは主張しない．
従って，位相的な Petersen star，標準 $M_{0,5}$ filling，実際の action lattices に適合する
integral affine chart は，それぞれ異なる条件として扱う必要がある．
\end{remark}

\section{作用周期からの integral affine 必要条件}
\label{trop:action-conditions-section}

ここでは存在定理を述べず，構成済み fibration の action coordinates を
unimodular な Petersen fan charts に合わせる際の必要条件を明示する．
正則 torus の $H_1(\TT^2;\ZZ)$ は integral cotangent lattice を与え，
primitive $\lambda$ による局所作用座標を
$I_\gamma=(2\pi)^{-1}\int_\gamma\lambda$ とする．
action coordinates の変換は整数線形変換と translation である．

\subsection{rational directions と annular areas}

一つの成分の比例曲率 collar で，$\beta_a=0$ の部分を考える．
normal circle の重みを $w>0$，connection を $\alpha$ とすると
\[
 \lambda=(\mu-w)\alpha,\qquad \dd\alpha=-\omega_P/w.
\]
normal cycle を $a$，regular level circle $\gamma$ の degree-one closed lift を $b$ とし，
\[
 h(\gamma)=\frac1{2\pi}\int_b\alpha
\]
と定める．$b$ の変更は $h$ に整数を加える．この基底で
\begin{equation}\label{trop:action-cone}
 (I_a,I_b)=(\mu-w)(1,h(\gamma)).
\end{equation}
Liouville cone の action coordinates は core に向かって一様に零へ収束するため，
core の像を原点に取れば，critical level に対応する放射辺の方向は $(1,h_*)$ である．
方向全体の符号は以下の条件に影響しない．

\begin{proposition}[ray と annulus の算術条件]
\label{trop:farey-condition}
上の collar に由来する core chart が，正則 faces 上の action lattice を保って
rational polyhedral fan に写るなら，各 critical value において
\begin{equation}\label{trop:rational-slope}
 h_*\in\mathbb Q
\end{equation}
でなければならない．

さらに，同じ成分内の二つの saddle levels の間の regular annulus $A$ が，
unimodular 二次元 cone の一つに対応するとする．両端の holonomies を，
$A$ 上で整合する lifts により $h_\pm=p_\pm/q_\pm$，$q_\pm>0$，
$\gcd(p_\pm,q_\pm)=1$ と書けば
\begin{equation}\label{trop:farey-determinant}
 |q_-p_+-q_+p_-|=1,
 \qquad
 \frac{\operatorname{area}(A)}{2\pi w}=\frac1{q_-q_+}
\end{equation}
が必要である．
\end{proposition}

\begin{proof}
ray に非零整数 vector が存在することは $h_*\in\mathbb Q$ と同値であり，
この性質は $\operatorname{GL}(2,\ZZ)$ で不変である．
$h_*=p/q$ のときの primitive direction は，符号を除いて $(q,p)$ となる．

$A$ の circle bundle は trivial なので，annular section を取って Stokes の定理を適用すると
\begin{equation}\label{trop:holonomy-area}
 h_+-h_-=-\frac1{2\pi w}\int_A\omega_P
\end{equation}
を得る．端点が critical leaves の場合は，regular annuli の exhaustion により同じ式を得る．
二本の primitive directions が cone の lattice basis になる条件は
\eqref{trop:farey-determinant} の第一式であり，第二式は
\eqref{trop:holonomy-area} によってこれと同値である．
\end{proof}

例えば $h_-=0$，$h_+=-2/3$ の directions はどちらも rational であるが，
primitive generators $(1,0),(3,-2)$ の determinant の絶対値は二である．
この period data は unimodular Petersen cone に適合しない．
また，葉層を境界に相対的に小さく isotopy すると，critical leaf が囲む area は連続的に変わる．
\eqref{trop:holonomy-area} により irrational な $h_*$ も生じるため，
位相的な marking を固定したことから rationality が自動的に従うわけではない．

\subsection{三つの sheets の saturated branch lattice}

三価の critical block で，向きを選んで
\[
 H_1(Q_v;\ZZ)=\ZZ a\oplus\ZZ b_1\oplus\ZZ b_2,
 \qquad b_3=b_1+b_2
\]
とする．三つの regular sheets の lattices は
$\langle a,b_1\rangle$，$\langle a,b_2\rangle$，$\langle a,b_3\rangle$ である．
critical level への period limits は $I_{b_3}=I_{b_1}+I_{b_2}$ を満たす．

\begin{proposition}[branch denominator の必要条件]
\label{trop:branch-denominator}
三つの sheets が共通の rank-one branch tangent lattice に沿って貼られ，
その lattice が各 sheet の action tangent lattice の saturated sublattice となると仮定する．
共通の circle action $I_a$ を三つの sheets で同じ関数とし，
branch の primitive tangent generator に対する $I_a$ の絶対勾配を $q>0$ とする．
このとき critical slopes は
\[
 h_1=p_1/q,\qquad h_2=p_2/q,\qquad h_3=(p_1+p_2)/q
\]
と書け，
\begin{equation}\label{trop:branch-coprime}
 \gcd(q,p_1)=\gcd(q,p_2)=\gcd(q,p_1+p_2)=1
\end{equation}
が必要である．特に $q$ は奇数である．
\end{proposition}

\begin{proof}
共通 primitive generator の三つの action-coordinate 表示を
$(q,p_1),(q,p_2),(q,p_3)$ とする．第一成分が同じ整数になるのは，$I_a$ が共通だからである．
period relation から $p_3=p_1+p_2$ となる．saturation は各表示の primitivity を意味するので
\eqref{trop:branch-coprime} を得る．$q$ が偶数なら $p_1,p_2$ はどちらも奇数であり，
$p_1+p_2$ は偶数となって第三の条件に反する．
\end{proof}

この条件と annulus の Farey 条件は両立すべき別の条件である．
例えば両端の $q$ が一なら，annulus を挟む numerator の差は一でなければならない．
抽象 fan の balancing identity \eqref{trop:petersen-balancing} だけでは，
これらの作用周期条件は検査されない．

\subsection{大域接着に必要な追加条件}

六次の Petersen choices に対して，以上は次のデータを確定する：
各 core の抽象 fan lattice，十五 cones の saturation と balancing，
実際の link fibration の common-circle determinants，および作用周期の必要条件である．
全体の action-compatible atlas を構成するには，さらに，各 face を通る
rank-two lattice identifications を選び，三価辺の common sublattices，
face transitions の cocycle condition，および実際の period data を同時に合わせる必要がある．

既存の focus--focus points の周囲には，適切な基底で
$\left(\begin{smallmatrix}1&1\\0&1\end{smallmatrix}\right)$ 型の非自明な monodromy がある
\cite{Evans,VuNgoc}．従って，action lattice と整合する通常の integral affine chart を
その点まで滑らかに延長することはできない．大域 atlas の定式化では，
focus points を除くか，その affine singularities を明示的に許す必要がある．
定義~\ref{atl:definition} はこの射程を明示し，定理~\ref{ga8:global} は
局所作用座標と実面積の接着によって，六次および五次の大域存在を証明する．

